% Français 9117, batch 6 — Gallica views 101–113, the last thirteen views of
% the volume.
% Pass: Opus 5 (claude-opus-5), 2026-09-14 — first pass, unchecked against
% the views by a human.
%
% What the views hold. The last ten written pages of the working manuscript
% of the « Considérations générales », in French prose, no mathematics; then
% the end of the volume. The hand is Germain's regular draft-for-the-press
% hand, struck and rewritten between the lines, some paragraphs opened by a
% square bracket in ink. Each leaf is written on both sides and the text runs
% on from page to page without a break.
%
%   v. 101     Opens a paragraph (« Dans cette manière de procéder, nos
%              auteurs suivent l'ordre établi par la réalité des
%              événemens ») — view 100 ends on a full stop (« dont elle ne
%              peut plus mesurer l'étendue. »), so nothing is cut at the
%              seam; view 101 begins with a capital and a bracket. The likeness of the arts to the facts that move us
%              rests on the identity of relations; each art has its module,
%              and the order of the parts of an action is not arbitrary.
%   v. 102–103 Variety of feeling prevents fatigue; music, of all the
%              languages of feeling, most demands a change of expression,
%              whence the improbable incidents of the lyric drama. The man
%              deaf to music finds the scenes ill-joined and thinks he has
%              judged what he cannot hear.
%   v. 104–105 The judgement of the amateur during the performance is the
%              same for all; the disputes of the schools vanish before the
%              execution, because uniformity of impression comes from the
%              faculties of feeling and diversity from opinion. The skilful
%              composer, like the great orators, never lets an impression
%              that is his triumph wear itself out: he replaces it by
%              another.
%   v. 106–107 What is extreme escapes attention; an extreme misfortune is
%              not at first felt more keenly than a lesser one, because
%              judgement needs to know its object; affections with an
%              outward cause tend to weaken. Letters and arts are our work
%              and their end our well-being: they must not carry their
%              action to the degree that troubles judgement. View 107 has in
%              its left margin a passage of some fifteen short lines struck
%              out line by line and crossed by a vertical stroke; nothing of
%              it can be read.
%   v. 108–109 Against the reproach that literature is exhausted: the ages of
%              hypotheses, when history and fable ran together, favoured the
%              imagination; creative power vanished with the credit of
%              fictions. The nations now are like a young man passing from
%              letters to serious studies; once the education is done each
%              thing recovers its true importance.
%   v. 110     Theories founded on incontestable truths will bring the
%              branches of knowledge to a harmony they once owed to
%              imagination alone; poets will show the universe as it is, and
%              genius will see that the art of creating was only that of
%              copying weak parts of a picture it will be given to paint
%              whole. The page ends with a full stop: it is the last written
%              page of the volume.
%   v. 111     A blank leaf, with the edge of v. 110 and the pencil
%              flourish on the guard at its left. Not transcribed.
%   v. 112–113 The back binding: marbled paper, the inside of the board and
%              the outside. Not transcribed. No microfilm target closes the
%              volume.
%
% The numbering on the leaves. Two series, both read on every view where
% they are given here, none inferred. (1) At the head of each page, top left,
% in ink, the page numbers of the manuscript, most of them crossed by a
% stroke as the earlier batches describe: 91 (underlined) v. 101, 92 v. 102,
% 93 v. 103, 94 v. 104, 95 v. 105, 96 v. 106, 97 v. 107, 98 v. 108, 99
% v. 109, 100 v. 110. They are recorded in a \note{} at each view. (2) At the
% top right of the rectos, in a thin, lighter hand, the foliation 50 v. 101,
% 51 v. 103, 52 v. 105, 53 v. 107, 54 v. 109 — the last agreeing with the
% catalogue's 54 leaves. These are written \folio{50r} … \folio{54r}, as in
% the other batches of the volume; nothing is written for the versos, which
% carry no folio. At the foot of the guard
% beside v. 110–111 there is a pencilled flourish, not a number.
%
% Other hands. Two names in another hand, as in the earlier batches, very
% likely the compositors' marks where each took up the copy (this pass's
% inference): in pencil at the head of v. 101, read uncertainly
% « Pennequin »; in the right margin of v. 106, in pale ink, « Derache »
% (batch 1 read the same name at its v. 19 as « Deraihe »). The Bibliothèque
% royale stamp on v. 110 is not transcribed.
%
% Conventions. Spelling of the page kept (« événemens », « appercevoir »,
% « connoître », « loix », « tems », « hazardées », « poëtes »,
% « parceque »); interlinear rewrites set in the line where they belong, the
% struck word before them, with a \note{} where the place matters. Small
% crosses and signs beside the punctuation (views 102, 105, 106, 107, 108,
% 110), apparently insertion or correction marks for the press, are not
% transcribed except where a note mentions one.
%
% What could not be read: the struck letters after « ma » (v. 101); a struck
% interlinear word over « devient », a struck interlinear word before « du
% mal », two struck lines round « première a eu une force réelle moins
% grande », and the struck marginal passage (all v. 107); a struck word
% before « indécise » (v. 108).
%
% Nothing here was checked against any published transcription.
\documentclass[11pt,a4paper]{article}
% The preamble shared by every transcription of the mathematicians' manuscripts in Gallica.
%
% It rests on one idea: **uncertainty compounds, it does not resolve.** A
% transcription that smooths over a doubtful word has destroyed the only thing
% it was made for — a reader must be able to tell, at every point, what was on
% the page and what was supplied. The apparatus macros below are therefore the
% critical apparatus, not decoration.
%
% The same commands are recognised by `scripts/render.mjs`, which produces the
% site's reading view, and by `scripts/tei.mjs`. A macro added here must be
% added there too, or rendering fails loudly — which is the wanted behaviour.
%
% Comments here are English, like the rest of the repository. The documents
% this preamble sets are French, and that is deliberate: see the skills.

\ProvidesPackage{germain}[2026/09/15 Transcriptions of mathematicians' manuscripts in Gallica]

\usepackage{amsmath,amssymb,amsthm}
\usepackage{mathrsfs}
% Kept from the parent preamble so the renderer's diagram support stays
% available; these manuscripts draw no commutative diagrams, and nothing here uses it.
\usepackage{tikz-cd}
\usepackage[english,french]{babel}
\usepackage{fontspec}

% --- Greek ------------------------------------------------------------------
% The text face stays fontspec's default, Latin Modern, which has no Greek:
% XeTeX drops every glyph it lacks with a line in the log and nothing on the
% page, and the Greek words the manuscripts quote vanished from the PDFs. So
% the Greek and Greek Extended blocks get a XeTeX character class of their own,
% and entering or leaving a run of it switches the family to CMU Serif — the
% same Computer Modern design, with polytonic Greek — and back. Only the
% family changes: a Greek word in \textit or \textbf keeps its shape and series.
% The transitions to and from every other class carry the tokens that class
% has towards an ordinary letter, so French spacing before « ; » or inside
% « » still holds after a Greek word. No grouping, so no brace or \endgroup
% can come out unbalanced; maths is left alone, since XeTeX inserts nothing
% there. Kept identical in `exercices.sty`.
\makeatletter
\newfontfamily\germainGreekfont{cmun}[
  Extension=.otf, NFSSFamily=germaingreek,
  UprightFont=*rm, ItalicFont=*ti, SlantedFont=*sl,
  BoldFont=*bx, BoldItalicFont=*bi, BoldSlantedFont=*bl]
\newXeTeXintercharclass\germain@greek
\def\germain@greekrange#1#2{%
  \@tempcnta=#1\relax
  \loop
    \XeTeXcharclass\@tempcnta=\germain@greek
    \ifnum\@tempcnta<#2\advance\@tempcnta\@ne\repeat}
\germain@greekrange{"0370}{"03FF}
\germain@greekrange{"1F00}{"1FFF}
\def\germain@greekprev{\rmdefault}
\def\germain@greekin{%
  \edef\germain@greekprev{\f@family}\fontfamily{germaingreek}\selectfont}
\def\germain@greekout{\fontfamily{\germain@greekprev}\selectfont}
\def\germain@greekjoin{%
  \XeTeXinterchartokenstate=\@ne
  \@tempcnta=\z@
  \loop
    \ifnum\@tempcnta=\germain@greek\else
      \XeTeXinterchartoks\germain@greek\@tempcnta=\expandafter{%
        \expandafter\germain@greekout\the\XeTeXinterchartoks\z@\@tempcnta}%
      \XeTeXinterchartoks\@tempcnta\germain@greek=\expandafter{%
        \the\XeTeXinterchartoks\@tempcnta\z@\germain@greekin}%
    \fi
    \ifnum\@tempcnta<4095\advance\@tempcnta\@ne\repeat}
% babel-french sets its own classes, and saves and restores the interchar
% state, each time a language is selected: join again after it does.
\addto\extrasfrench{\germain@greekjoin}
\addto\extrasenglish{\germain@greekjoin}
\AtBeginDocument{\germain@greekjoin}
\makeatother

\usepackage{geometry}
\geometry{a4paper,margin=25mm,marginparwidth=28mm}
\usepackage{marginnote}
\usepackage{xcolor}
\usepackage[normalem]{ulem}
\setlength{\emergencystretch}{3em}
\usepackage{hyperref}
\hypersetup{colorlinks=true,linkcolor=black,urlcolor=black,pdfborder={0 0 0}}

% --- Batch metadata ---------------------------------------------------------
% Taken as they stand by the rendering script, which shows them at the head of
% the reading view. They fix what the file claims to cover. The macro names are
% the parent project's, so that the scripts stay shared; \folder here means the
% volume's slug — `fr-9115` — and \pages counts Gallica views.

\newcommand{\folder}[1]{\gdef\grothFolder{#1}}
\newcommand{\batch}[1]{\gdef\grothBatch{#1}}
\newcommand{\pages}[2]{\gdef\grothFirst{#1}\gdef\grothLast{#2}}
\newcommand{\foldertitle}[1]{\gdef\grothFoldertitle{#1}}

% The holder's own shelfmark, printed as they write it: « Français 9115 ».
\newcommand{\shelfmark}[1]{\gdef\grothShelfmark{#1}}

% The dating, copied verbatim from the holder's catalogue — never inferred
% here. « XIXe siècle » is the BnF's; a pass that dates a leaf from its content
% says so in a \note{}, as its own inference.
\newcommand{\dating}[1]{\gdef\grothDating{#1}}

% The status notice, printed in the title block and repeated as a diagonal
% watermark on every page of the PDF. The manuscripts are in the public
% domain; what this line declares is not a right but a quality — an
% unverified first pass by a machine. The skills write the line; a file
% without it gets no watermark, so leaving it out is visible in the output.
\newcommand{\watermark}[1]{\gdef\grothWatermark{#1}}

\gdef\grothFolder{?}\gdef\grothBatch{}\gdef\grothFirst{?}\gdef\grothLast{?}%
\gdef\grothFoldertitle{}\gdef\grothShelfmark{}\gdef\grothDating{}\gdef\grothWatermark{}

% --- The résumé -------------------------------------------------------------
\newenvironment{resume}
  {\small\setlength{\parskip}{0.45em}}
  {\par\medskip\hrule\medskip}

% --- Keywords ---------------------------------------------------------------
% In English, closing the modernised reading's résumé: the modern vocabulary
% under which to search for what the leaves construct. The manifest extracts
% them to tag the volume — this line is the single source of the tags.
\newcommand{\keywords}[1]{\par\smallskip\noindent
  {\footnotesize\textsc{Keywords} --- \textit{#1}}\par}

% --- The critical apparatus -------------------------------------------------

% The Gallica view number, in the margin. This is the anchor that lets a
% disagreement over a formula be settled by looking at *one* image rather than
% a volume of 750. The site uses it to turn the facsimile as the transcription
% is scrolled. A view is one scanned image: usually one side of a leaf.
\newcommand{\page}[1]{%
  \marginnote{\footnotesize\color{gray!70}\textsf{v.\,#1}}%
  \label{page:#1}\ignorespaces}

% The BnF's pencilled foliation, where the leaf carries it and a pass can read
% it: « 198r ». This is what the literature cites — « ff. 198r–208v » — and
% the one line that turns a citation into a view. Recorded, never inferred:
% a folio a pass cannot read is not written.
\newcommand{\folio}[1]{%
  \marginnote{\footnotesize\color{gray!50}\textsf{f.\,#1}}\ignorespaces}

% The modernised reading's coarser counterpart to \page{}: which views a
% section is drawn from.
\newcommand{\pagerange}[2]{%
  \marginnote{\footnotesize\color{gray!70}\textsf{v.\,#1--#2}}%
  \ignorespaces}

% Illegible: never guessed.
\newcommand{\ill}{\textcolor{red!60!black}{[\,\ldots\,]}}

% A doubtful reading: the word is given, and flagged as uncertain.
\newcommand{\uncertain}[1]{\uline{#1}}

% An editorial addition — a missing letter, an implied word.
\newcommand{\add}[1]{\textcolor{blue!50!black}{[#1]}}

% Struck out by the writer. What was crossed out is part of the document.
\newcommand{\struck}[1]{\sout{#1}}

% The transcriber's note, on the state of the page or the mathematics.
\newcommand{\note}[1]{\par\smallskip{\small\color{gray!60!black}\textit{[#1]}}\par\smallskip}

% A marginal note of hers, distinct from the body of the text.
\newcommand{\marginal}[1]{\par\smallskip\noindent\hspace{1em}%
  {\small\color{gray!60!black}#1}\par\smallskip}

% --- Title ------------------------------------------------------------------
% The title block is French, like the documents themselves.

\AddToHook{shipout/background}{%
  \ifx\grothWatermark\empty\else
    \begin{tikzpicture}[remember picture, overlay]
      \node[rotate=52, scale=3.4, align=center, text=gray!18,
            font=\sffamily\bfseries]
        at (current page.center) {\grothWatermark};
    \end{tikzpicture}%
  \fi}

\AtBeginDocument{%
  \begin{center}
    {\large\bfseries Manuscrits de mathématiciens (Gallica) ---
      \ifx\grothShelfmark\empty\grothFolder\else\grothShelfmark\fi}\\[2pt]
    {\itshape\grothFoldertitle}\\[4pt]
    {\small vues \grothFirst--\grothLast
      \ifx\grothBatch\empty\else\ (lot \grothBatch)\fi}\\[2pt]
    \ifx\grothDating\empty\else
      {\small\color{gray!60!black}Datation du catalogue : \grothDating}\\[2pt]
    \fi
    \ifx\grothWatermark\empty\else
      {\small\bfseries\color{red!45!gray}\grothWatermark}\\[2pt]
    \fi
    {\footnotesize\color{gray!60!black}Transcription établie d'après les images IIIF de Gallica.
     Source gallica.bnf.fr / Bibliothèque nationale de France.\\
     Les lectures incertaines sont soulignées, les illisibles notées [\,\ldots\,],
     les ajouts entre crochets ; \textsf{v.} numérote les vues de Gallica,
     \textsf{f.} la foliotation au crayon de la BnF.}
  \end{center}
  \vspace{1em}}

\endinput


\folder{fr-9117}
\shelfmark{Français 9117}
\batch{6}
\pages{101}{113}
\dating{1801-1900}
\watermark{Lecture automatique\\non vérifiée}
\foldertitle{« Considérations générales sur l'état des sciences et des lettres aux différentes époques de leur culture, » par Sophie GERMAIN.}

\begin{document}

\page{101}\folio{50r}
\note{en tête de la page, à l'encre et soulignée : « 91 » ; au crayon, d'une autre main, un nom : \uncertain{Pennequin}}

\note{un crochet ouvrant à l'encre devant le premier mot}
Dans cette \struck{ma\ill{}} manière de procéder, nos auteurs suivent
l'ordre établi par la réalité des événemens qui peuvent affecter notre
sensibilité. Il arrive, en effet, \struck{que d} tous les jours qu'un détail
nouveau, une circonstance imprévue qui accompagne un malheur déjà connu en
redoublent l'impression, au point de nous \uncertain{jetter} dans le désespoir.

La ressemblance de nos arts \uncertain{entr'eux} et avec la nature des faits qui
nous émeuvent tient à l'identité de rapports, sans laquelle il n'existeroit
pour nous aucun sentiment vrai, aucune idée claire. chaque art, aussi bien que
la réalité des événemens dont ils imitent les impressions, \struck{ont tout}
son module particulier ; et c'est en quoi ils diffèrent les uns des autres ;
mais, ce module étant adopté, rien n'est plus arbitraire entre les diverses
parties de l'action. leur liaison est tellement nécessaire que, si elle étoit
intervertie, nous n'appercevrions plus aucune suite d'action, nous
n'éprouverions plus aucune gradation d'intérêt.

Lorsque les différentes parties d'une composition ont été coordonnées avec
art, l'âme se plaît à en parcourir le développement :

\page{102}
\note{en tête de la page, à l'encre : « 92 »}

la variété des sentimens prévient la fatigue, et \struck{\uncertain{répand l'a}}
soutient l'attention. Plus la composition a de force et d'énergie, plus aussi,
il \struck{devient} est nécessaire de ne pas s'arrêter trop longtems
\struck{à une situation o} à l'expression d'une seule et même impression.
Celle de la tristesse, par exemple deviendroit assoupissante ; et les accens du
désespoir, s'ils étoient trop multipliés, finiroient par n'être plus entendus.

La musique, pour qui entend son langage, est de tous les genres de discours qui
ont les sentimens pour objets, celui qui exige le plus \struck{impérieusement}
la variété du \uncertain{stile} ; parcequ'il est aussi celui dont les
impressions sont le plus pénétrantes.

Cette nécessité du changement d'expression est tellement impérieuse dans l'art
musical qu'on est obligé, pour y satisfaire, d'introduire dans le drame lyrique,
les incidens les plus invraisemblables, lorsqu'il ne s'en présente aucun, dans
la suite naturelle de l'action, qui puisse fournir au compositeur,
\struck{l'occasion de varier} des sentimens de genres différens à mettre en
scène.
\note{sous la virgule qui suit « compositeur », une petite croix}

La partie poëtique d'un tel drame expose nécessairement une action

\page{103}\folio{51r}
\note{en tête de la page, à l'encre : « 93 »}

déterminée ; tandis que sa partie musicale \struck{de ce même drame} saisit
l'occasion qui lui est offerte, pour exprimer des sentimens généraux et
abstraits qu'amène le récit du poëte ; \struck{amène en scène}.
\note{« amène », dans « qu'amène », est écrit serré entre « qu' » et « le »}

C'est de cette diversité d'actions que résulte celle du jugement des
spectateurs.
\note{un crochet ouvrant à l'encre devant « S'il arrive »}
S'il arrive qu'un homme peu sensible au langage musical, s'avise de vouloir
aller entendre quelqu'un des chefs d'œuvres dramatiques qui excitent
l'enthousiasme des amateurs, \struck{ce spectateur} \struck{intrig} il cherche,
comme malgré lui, une distraction dans le sujet particulier de la pièce, contre
les accords dont il ne sent pas le mérite. Il s'étonne du peu de liaison des
scènes ; il ne voit pas avec quel art on a su ménager ces transitions
nécessaires entre des genres de beautés tellement énergiques que
\struck{leur impress} l'impression de chacune d'elles seroit une horrible
fatigue si elle étoit prolongée. Il sort en croyant avoir jugé ce qu'il ne
sait pas entendre ; et croit avoir fait une plaisanterie mordante en nous
disant qu'il semble que le bon sens doit être partout à sa place.

\page{104}
\note{en tête de la page, à l'encre : « 94 », écrit dans l'angle d'un crochet ouvrant tracé devant le premier mot}

L'amateur est pleinement satisfait ; il n'a pas même remarqué les défauts
prétendus que l'on semble trouver dans une composition dont il admire toutes
les parties.

\note{un trait horizontal suit « parties. » en fin de ligne}
Un tel jugement est univoque. L'âge, le sexe, la position sociale, les
connoissances les plus approfondies dans des sujets différens, ou leur défaut
total, \uncertain{n'établissent} aucun dissentiment sur l'impression
\struck{récente} actuelle ou récente, produite par un bon ouvrage. Toutes ces
personnes, si dissemblables d'ailleurs dans leurs goûts et leurs habitudes, se
réunissent dans une opinion qui leur est commune. A la vérité, cet accord
universel se trouble bientôt. — ce n'est qu'alors que l'impression dure encore
qu'elle dicte à chacun des jugemens semblables. On peut même remarquer une
chose qui paroîtroit inexplicable à tout homme étranger aux impressions
musicales, c'est que les discussions quelquefois très vives sur le mérite des
différentes écoles, sur celui de tel ou tel exécutant, disparoissent, ou ne
sont plus

\page{105}\folio{52r}
\note{en tête de la page, à l'encre : « 95 »}

soutenues qu'avec peine et par \uncertain{pur} entêtement, en présence de
l'exécution. C'est que cette diversité est \struck{une} dans l'opinion ;
tandis que l'uniformité d'impression \struck{est} naît de celle des facultés de
sentir. La première est une force constante ; l'autre a besoin d'être mise en
action pour acquérir \struck{une} de la valeur ; le souvenir la reproduit
imparfaitement, et son effet, s'affoiblissant en raison de l'éloignement de la
cause qui l'a produit, laisse à la première \struck{toute} une prépondérance
marquée.
\note{« de l'éloignement » est ajouté dans l'interligne, avec un signe de renvoi ; au-dessus de « produit, laisse à », deux mots ajoutés dans l'interligne : \uncertain{finit par}}

\note{un fort trait horizontal, à gauche, sous « la première »}
C'est par une raison toute semblable que l'amateur, quelque \uncertain{sûrs}
d'ailleurs la puissance de sa raison et la délicatesse de son goût littéraire,
n'est point choqué des invraisemblances et du peu de liaison entre les
\struck{des} scènes lyriques. \struck{par une raison toute semblable}.
\note{« C'est par une raison toute semblable » est écrit dans l'interligne ; la même locution, biffée, suit « scènes lyriques » ; « entre les » est ajouté dans l'interligne au-dessus de « des », biffé}
Pénétré d'une vive émotion \struck{ému} par ce qu'il entend actuellement, il
n'a pas le loisir d'appercevoir quel en est l'à-propos. Le compositeur habile,
semblable en cela aux grands orateurs s'empare de l'attention de ses
auditeurs ; leur âme est pour ainsi dire dans sa main ; il dispose de leurs
sentimens. Après avoir porté le trouble dans leurs facultés, et, lorsqu'il sent
que ce trouble ne peut plus croître, il ne veut pas, \struck{\uncertain{le}}
laisser affoiblir une impression qui est son triomphe ; il \struck{cherche une}
veut qu'une impression différente vienne remplacer la première,
\note{« ému », biffé, est dans l'interligne au-dessus de « par »}

\page{106}
\note{en tête de la page, à l'encre : « 96 »}

et conserve ainsi à l'empire qu'il exerce sur nos âmes toute sa force et toute
sa puissance. — Avec quel art il sait en ménager les moyens ! Il connoît,
\struck{jusqu'à quel} il mesure, les impressions qu'il fait naître ; il sait
combien de tems nous pouvons en être possédés ; il se garde d'atteindre la
limite de nos facultés. Ce qui est extrême échappe à notre attention ; la
faculté de sentir a ses bornes. Et cette remarque, vraie à l'égard des
impressions dont les arts disposent, tient à la nature de notre être. Elle est
vraie en tout genre. Un extrême malheur, dans les premiers momens de son
apparition n'est pas senti plus vivement qu'un malheur moindre ; mais les
instans qui succèdent nous font reconnoître les différences. \struck{parceq}
Et cela par une raison applicable aussi à toute chose, parcequ'elle est vraie,
et que la vérité est universelle, ou, en d'autres termes \struck{ce qui est la
même chose} parceque les loix de l'être sont partout les mêmes.
\note{« en d'autres termes » est écrit dans l'interligne, au-dessus de « ce qui est la même chose », biffé}

\note{un crochet ouvrant à l'encre devant « Cette » ; dans la marge de droite, d'une autre main et d'une encre plus pâle : « Derache »}
Cette \struck{remarque} raison, la voici. nos jugemens \struck{ne sont tant}
pour être éclairés ont besoin de la connoissance parfaite des choses qui en sont

\page{107}\folio{53r}
\note{en tête de la page, à l'encre : « 97 »}

l'objet ; une impression d'une extrême force \struck{porte} nous jette dans le
trouble ; \struck{Dans} notre âme \struck{elle} devient incapable de comparer,
et par conséquent de juger, la force, \struck{et, pour ainsi dire}
l'intensité du malheur qu'elle éprouve.
\note{« nous jette dans » est écrit dans l'interligne au-dessus de « porte », biffé}

Mais il est dans la nature des affections qui ont une cause extérieure de
tendre à s'affoiblir. \struck{Un} Le chagrin cuisant \struck{devient}
\struck{\ill{}} semble devenir de plus en plus poignant, parceque, l'impression,
qui d'abord avoit été assez forte pour nous empêcher d'en mesurer l'étendue
\struck{\ill{}} du mal, diminuant réellement, laisse à la réflexion le pouvoir
de nous en montrer les différentes faces. Si la cause \struck{\ill{}} première
a eu une force réelle moins grande \struck{\ill{}} son effet aura été connu
plutôt ; et l'affoiblissement de l'impression, loin de tendre à augmenter notre
chagrin, devra, au contraire, amener un soulagement sensible à nos peines.
\note{« Le » est écrit au-dessus de « Un », biffé ; au-dessus de « devient », biffé, un mot de l'interligne, biffé lui aussi ; « d'en » est récrit dans l'interligne au-dessus du même mot biffé ; « du mal » est écrit sous la ligne, après un mot de l'interligne biffé et illisible ; une ligne entière de l'interligne au-dessus de « première a eu une force réelle moins grande » et la ligne qui la suit sont biffées et illisibles}
\marginal{\struck{\ill{}} — dans la marge de gauche, un passage d'une quinzaine de lignes courtes, raturé ligne par ligne et barré d'un trait vertical, en regard de « il est dans la nature des affections » jusqu'à « laisse à la réflexion »}

Les lettres, les arts, \struck{tout fait par nous et pour} sont notre ouvrage ;
leur but est notre bien-être ; la première règle qu'ils doivent suivre est donc
de ne pas porter leur action jusqu'à ce degré extrême qui trouble notre
jugement ; ou s'il est dans leur nature de savoir l'atteindre, et que de grandes
beautés \struck{doivent} puissent résulter de ce genre d'impressions, ils
doivent se hâter de nous tirer d'une position pénible, \struck{elle} nous
\note{« sont notre ouvrage ; leur but » et « est » sont écrits dans l'interligne ; « donc » et « puissent » aussi ; « ou » est récrit dans l'interligne au-dessus du même mot biffé}

\page{108}
\note{en tête de la page, à l'encre : « 98 »}

transportant notre attention dans une région moins sombre.

On reproche à la littérature son épuisement, et au goût de l'exactitude sa
sécheresse. Il semble que l'imagination ait perdu sa puissance lorsque la
raison établit son empire. Nous reconnoîtrons que \struck{le passage entre} les
tems où des hypothèses, plus ou moins heureuses, formoient les richesses
intellectuelles ; où l'homme, par conséquent au lieu de chercher l'appui des
vérités particulières, ou, ce qui est la même chose celui de la réalité des
faits, trouvoit en lui-même les convenances auxquelles il assujettissoit la
nature entière, étoient \struck{plus} très favorables au développement de
l'imagination. Alors une foule de doctrines hazardées, ne contrastoient pas avec
ce qu'on nommoit la science. les fictions poëtiques étoient revêtues d'un charme,
qu'elles ne devoient pas seulement à leur grâce ; une \uncertain{demi} croyance
\struck{\ill{}} \struck{indécise} permettoit d'admettre qu'elles pouvoient avoir
eu une sorte de réalité. L'histoire et la fable, se confondoient dans leurs
limites ; ce que les uns regardoient comme de simples allégories étoient, aux
yeux des autres, le récit de faits merveilleux
\note{« très » est écrit dans l'interligne au-dessus de « plus », biffé}

\page{109}\folio{54r}
\note{en tête de la page, à l'encre : « 99 »}

Cette disposition des esprits donnoit sans doute à l'art de bien dire une
importance qu'il ne peut conserver au même degré lorsque la principale condition
à remplir est celle de dire vrai. — La faculté créatrice a disparu avec le
crédit des fictions. Mais, s'il est dans le \struck{nature de notre manière}
caractère de notre culture intellectuelle d'attacher plus d'importance à la
solidité des doctrines qu'à ce qu'elles pourroient offrir de brillant ; si nous
voulons que la raison domine toutes les productions de l'esprit ; si même nous
sentons le goût des recherches attiédir notre imagination ; ne désespérons pas
d'arriver à une époque plus heureuse où nous saurons unir toutes nos facultés
dans des productions d'un genre nouveau.
\note{« à remplir » et « caractère » sont écrits dans l'interligne ; « le » est récrit sur « la »}

Les nations éprouvent aujourd'hui l'impression que recevroit un jeune homme qui,
après \struck{avoir} s'être longtems occupé de \struck{l'ancienne} littérature
se trouveroit, porté, par le cours de ses études, vers les connoissances
sérieuses. Le charme de ses premières occupations l'abandonneroit dans cette
\struck{route} nouvelle ; une curiosité vive en \struck{occuperoit} prendroit la
place ; mais, après l'éducation achevée, chaque chose \struck{reprendroit}
recouvreroit à ses yeux l'importance qui lui appartient réellement.
\note{« prendroit » et « recouvreroit » sont écrits dans l'interligne}

Nous arriverons à une époque semblable. Et, comme l'éducation des sociétés
consiste moins à rendre vulgaires les connoissances déjà anciennes qu'à en
acquérir de nouvelles ; comme nous marchons à grands pas vers la création de
théories

\page{110}
\note{en tête de la page, à l'encre : « 100 »}

fondées sur des vérités incontestables, nous devons finir par amener les
différentes branches de notre savoir à une harmonie, qu'elles durent autrefois
à notre seule imagination.

Tant de vérités, de genres différens, \struck{pour ainsi dire} groupées autour
d'une vérité première, qui est le fait principal du sujet, mettront dans tout
son jour l'identité de rapports, entre le module de chaque science, de chaque
art et les diverses parties de cette science ou de cet art. Les loix de l'être,
les conditions du vrai, ainsi présentées à la fois sous mille faces
différentes, échaufferont alors l'imagination ; un enthousiasme nouveau fondé
sur une base plus solide, \struck{saura nous présenter le charme de cet unité}
que celui qui sut embellir d'heureuses fictions, inspirera nos poëtes et nos
orateurs. Au lieu de créer l'univers suivant les caprices de nos volontés, ils
nous le montreront tel qu'il est réellement ; et, si jamais le génie entre dans
cette route nouvelle, il verra avec admiration que l'art de créer, n'a été que
celui de copier et de transporter en d'autres lieux \struck{les parties}
de foibles parties d'un tableau qu'il lui sera donné de savoir peindre dans son
entier.
\note{« savoir » est taché d'encre ; « à la fois » est écrit dans l'interligne ; la page s'achève ici sur un point, et le texte ne continue pas au-delà dans le volume}

\end{document}
